\documentclass[10pt,a4paper]{article}
\usepackage[T1]{fontenc}
\usepackage{lmodern,microtype}
\usepackage{amsmath,amssymb,mathtools,mathpartir}
\usepackage[margin=21mm,headheight=14pt]{geometry}
\usepackage{booktabs,listings,fancyhdr,xcolor}
\usepackage[hidelinks]{hyperref}
\definecolor{ink}{HTML}{263849}
\pagestyle{fancy}
\fancyhf{}
\fancyhead[L]{\small A Minimal Calculus of Implicit Implementations}
\fancyhead[R]{\small Formal specification}
\fancyfoot[C]{\thepage}
\renewcommand{\headrulewidth}{0.3pt}
\setlength{\parindent}{0pt}
\setlength{\parskip}{5pt}
\lstset{basicstyle=\ttfamily\small,columns=fullflexible,keepspaces=true,
  frame=single,framerule=0.2pt,breaklines=true,showstringspaces=false}
\newcommand{\Int}{\mathsf{Int}}
\newcommand{\Str}{\mathsf{String}}
\newcommand{\Bool}{\mathsf{Bool}}
\newcommand{\TS}{\mathsf{ToString}}
\newcommand{\wf}[1]{\Gamma\vdash #1\;\mathsf{type}}
\newcommand{\E}[3]{\Gamma;\Delta\vdash #1:#2\rightsquigarrow #3}
\newcommand{\R}[2]{\Gamma;\Delta\vdash #1\Longrightarrow #2}
\newcommand{\Pj}[2]{\Gamma;\Delta\vdash_{\mathrm{ev}} #1:#2}
\newcommand{\need}[1]{\mathsf{need}[#1]}
\newcommand{\hole}[1]{?_{#1}}
\newcommand{\ilam}[3]{\lambda\{#1:#2\}.\,#3}
\newcommand{\lam}[3]{\lambda(#1:#2).\,#3}
\newcommand{\letin}[3]{\mathsf{let}\;#1=#2\;\mathsf{in}\;#3}
\newcommand{\impl}[4]{\mathsf{impl}\;#1:#2=#3\;\mathsf{in}\;#4}
\newcommand{\argop}[1]{\operatorname{arg}(#1)}
\newcommand{\resop}[1]{\operatorname{res}(#1)}
\newcommand{\erase}[1]{\lvert #1\rvert}
\newcommand{\eval}[3]{#1\vdash #2\Downarrow #3}
\newcommand{\kw}[1]{\mathsf{#1}}
\newcommand{\fields}[1]{\{#1\}}
\begin{document}
{\Large\bfseries\color{ink} A Minimal Calculus of Implicit Implementations}\par
{\small Scala 3 reference model \hfill 28 October 2026}

\section{Types, expressions, and contexts}
\[
\begin{array}{rcl}
A,B,P,Q&::=&\alpha\mid\Int\mid\Bool\mid\Str\mid A\to B\mid A\Rightarrow B
 \mid\forall\alpha.A\mid\{\ell_j:A_j\}_{j\in J}\\[2pt]
\Gamma&::=&\cdot\mid\Gamma,\alpha\mid\Gamma,x:A\\
\Delta&::=&[\,]\mid(d:A)::\Delta
\end{array}
\]
$\Gamma$ binds type variables and ordinary values; $\Delta$ binds available implementation
values, nearest first. $P,Q$ are ordinary types, not a distinct predicate sort.
Record labels are distinct and canonically ordered. Equality is structural, modulo
$\alpha$-equivalence and lexical type-synonym expansion. There is no subtyping.

\textbf{First expression AST} (after type-name expansion):
\[
\begin{array}{rcl}
e&::=&x\mid d\mid n\mid b\mid s\mid\lam{x}{A}{e}\mid\ilam{d}{A}{e}
 \mid e(e)\mid e\{e\}\\
&&\mid\Lambda\alpha.e\mid e[A]\mid\letin{x}{e}{e}\mid\impl{d}{A}{e}{e}\\
&&\mid\{\ell_j=e_j\}_{j\in J}\mid e.\ell\mid(e:A)\mid\need A
 \mid e_1\mathbin{\odot}e_2
\end{array}
\]
The actual named AST additionally has $\kw{type}\ C[\bar\alpha]=A\ \kw{in}\ e$.
Both $x$ and $d$ are parsed as names; lexical name resolution distinguishes their binders.
In rules, bound names are globally fresh by $\alpha$-renaming. This convention preserves
ordinary source shadowing without conflating the two environments.

\textbf{Unresolved core}:
\[
\begin{array}{rcl}
c&::=&x\mid d\mid n\mid b\mid s\mid\lam{x}{A}{c}\mid\ilam{d}{A}{c}
 \mid c\,c\mid c\{c\}\\
&&\mid\Lambda\alpha.c\mid c[A]\mid\letin{x:A}{c}{c}\mid\impl{d}{A}{c}{c}\\
&&\mid\{\ell_j=c_j\}_{j\in J}\mid c.\ell\mid\hole A
 \mid c_1\mathbin{\odot}c_2
\end{array}
\]
The core has no aliases, traits, definitions, modules, or method-resolution operation.
$\hole A$ is a deferred request, not a runtime term.

\textbf{Primitive binary operators}:
\[
\begin{array}{c|c|c}
\odot & \argop\odot & \resop\odot\\\hline
+,-,\times & \Int & \Int\\
\mathsf{==},<,> & \Int & \Bool\\
\mathbin{+\!+} & \Str & \Str
\end{array}
\]
\[
\begin{aligned}
\delta_+(m,n)&=m+n, &\delta_-(m,n)&=m-n, &\delta_\times(m,n)&=mn,\\
\delta_{==}(m,n)&=(m=n),&\delta_<(m,n)&=(m<n),&\delta_>(m,n)&=(m>n),\\
\delta_{++}(s,t)&=st.&&&&
\end{aligned}
\]
Here $b\in\{\mathsf{true},\mathsf{false}\}$; comparisons produce Boolean values.

\textbf{Type formation} (contexts are assumed well formed):
\begin{mathpar}

\inferrule*[right=W-Var]{\alpha\in\Gamma}{\wf\alpha}

\inferrule*[right=W-Base]{B\in\{\Int,\Bool,\Str\}}{\wf B}

\inferrule*[right=W-Arrow]{\wf A\\\wf B\\\star\in\{\to,\Rightarrow\}}
 {\wf{A\star B}}

\inferrule*[right=W-All]{\Gamma,\alpha\vdash A\;\kw{type}\\\alpha\ \kw{fresh}}
 {\wf{\forall\alpha.A}}

\inferrule*[right=W-Record]{\wf{A_j}\quad(j\in J)\\\ell_j\text{ distinct}}
 {\wf{\{\ell_j:A_j\}_{j\in J}}}
\end{mathpar}

\newpage
\section{Practical syntax and the first AST}
\begin{lstlisting}
trait ToString[A] = { toString: A -> String }

impl ToString[String] {
  toString = (text: String) => text
}

def display[A where ToString[A]](x: A): String = x.toString

display[String]("Ada")
\end{lstlisting}
A declaration sequence is a nested expression, not a module. The parser constructs the
right-hand sides below directly; $b$ is the remainder of the sequence.
\[
\begin{aligned}
&\kw{trait}\ C[\bar\alpha]=A;\ b
 &&\equiv\kw{type}\ C[\bar\alpha]=A\ \kw{in}\ b,\\[2pt]
&\kw{type}\ C[\bar\alpha]=A;\ b
 &&\equiv\kw{type}\ C[\bar\alpha]=A\ \kw{in}\ b,\\[2pt]
&\kw{val}\ x=e;\ b
 &&\equiv\letin{x}{e}{b},\\[2pt]
&\kw{val}\ x:A=e;\ b
 &&\equiv\letin{x}{(e:A)}{b},\\[2pt]
&\kw{impl}\ d:A=e;\ b
 &&\equiv\impl{d}{A}{e}{b}.
\end{aligned}
\]
For $\bar P=P_1,\ldots,P_k$ and fresh $d_1,\ldots,d_k$:
\[
\begin{aligned}
&\kw{def}\ f[\bar\alpha\ \kw{where}\ \bar P](\bar x:\bar A):B=e;\ b\\
&\quad\equiv
 \letin{f}{\Lambda\bar\alpha.\,
 \lambda\{d_1:P_1\}.\cdots\lambda\{d_k:P_k\}.
 \lambda(\bar x:\bar A).(e:B)}{b}.
\end{aligned}
\]
\[
\begin{aligned}
&\kw{impl}\ i[\bar\alpha\ \kw{where}\ \bar P]:Q\ \{\bar\ell=\bar e\};\ b\\
&\quad\equiv\impl{i}{\forall\bar\alpha.P_1\Rightarrow\cdots\Rightarrow P_k\Rightarrow Q}
 {\Lambda\bar\alpha.\lambda\{d_1:P_1\}.\cdots\lambda\{d_k:P_k\}.
 \{\bar\ell=\bar e\}}{b}.
\end{aligned}
\]
An implementation body after \texttt{=} may be any expression. A record body without
\texttt{=} is the same expression. Anonymous \texttt{impl Q \{...\}} introduces a fresh,
unspellable binding name. Ordinary parameter lists are curried. The parser uses precedence
$\mathsf{postfix}>\times>\{+,-\}>\mathbin{+\!+}>\{==,<,>\}$, with
left-associative binary operators.

\subsection*{Lexical type aliases and qualified references}
A nonrecursive synonym $C[\bar\alpha]=B$ expands
$C[\bar A]$ to $B[\bar A/\bar\alpha]$ in its lexical body. Expansion is
capture avoiding: $B$ retains its declaration-time type-variable bindings.
$|\bar A|=|\bar\alpha|$ is required. A term binding shadows an equally named synonym.
If a name resolves to a synonym rather than a term binding:
\[
 C[\bar A]\text{ in term position}\quad\equiv\quad\need{C[\bar A]}.
\]
Thus \texttt{ToString[A].toString(x)} denotes a projection from a deferred
implementation request. The core type is directly
\[
\TS[A]\triangleq\{\kw{toString}:A\to\Str\}.
\]
\subsection*{Judgments and implementation representation}
\[
\boxed{\E e A c}\qquad
\boxed{\R c t}\qquad
\boxed{H\vdash t:T}\qquad
\boxed{\eval\rho t v}
\]
These are, respectively, name/type elaboration, evidence resolution, explicit-target
typing, and evaluation. The implementation uses separate de Bruijn indices for $\Gamma$
and $\Delta$, and a type-variable depth. The final target has one term-index space.

\newpage
\section{Typing and elaboration: functions and bindings}
All annotation types must satisfy $\wf A$. Instance selection does not occur in this judgment.
\begin{mathpar}

\inferrule*[right=E-Var]{x:A\in\Gamma}{\E x A x}

\inferrule*[right=E-Instance]{d:A\in\Delta}{\E d A d}


\inferrule*[right=E-Lambda]{\Gamma,x:A;\Delta\vdash e:B\rightsquigarrow c}
 {\E{\lam{x}{A}{e}}{A\to B}{\lam{x}{A}{c}}}


\inferrule*[right=E-Implicit-Lambda]{\Gamma;(d:A)::\Delta\vdash e:B\rightsquigarrow c}
 {\E{\ilam{d}{A}{e}}{A\Rightarrow B}{\ilam{d}{A}{c}}}


\inferrule*[right=E-Type-Lambda]{\Gamma,\alpha;\Delta\vdash e:A\rightsquigarrow c\\
 \alpha\notin\mathrm{FV}(\Gamma,\Delta)}
 {\E{\Lambda\alpha.e}{\forall\alpha.A}{\Lambda\alpha.c}}


\inferrule*[right=E-Type-App]{\E e {\forall\alpha.A} c\\\wf B}
 {\E{e[B]}{A[B/\alpha]}{c[B]}}


\inferrule*[right=E-Implicit-App]{\E{e_1}{A\Rightarrow B}{c_1}\\\E{e_2}{A}{c_2}}
 {\E{e_1\{e_2\}}{B}{c_1\{c_2\}}}


\inferrule*[right=E-App]{\E{e_1}{T}{c_1}\\c_1:T\triangleright c'_1:A\to B\\\E{e_2}{A}{c_2}}
 {\E{e_1(e_2)}{B}{c'_1\,c_2}}


\inferrule*[right=E-Let]{\E{e_1}{A}{c_1}\\\Gamma,x:A;\Delta\vdash e_2:B\rightsquigarrow c_2}
 {\E{\letin{x}{e_1}{e_2}}{B}{\letin{x:A}{c_1}{c_2}}}


\inferrule*[right=E-Impl]{\E{e_1}{A}{c_1}\\\Gamma;(d:A)::\Delta\vdash e_2:B\rightsquigarrow c_2}
 {\E{\impl{d}{A}{e_1}{e_2}}{B}{\impl{d}{A}{c_1}{c_2}}}


\inferrule*[right=E-Ascribe]{\E e A c}{\E{(e:A)}{A}{c}}

\inferrule*[right=E-Need]{\wf A}{\E{\need A}{A}{\hole A}}
\end{mathpar}

\subsection*{Expose an ordinary application}
\begin{mathpar}

\inferrule*[right=Expose-Stop]{T\ne(P\Rightarrow Q)\text{ for every }P,Q}
 {c:T\triangleright c:T}

\inferrule*[right=Expose-Step]{c\{\hole P\}:Q\triangleright c':T}
 {c:P\Rightarrow Q\triangleright c':T}
\end{mathpar}
Only ordinary application invokes this relation. In particular, returning an implicit
function, referring to it by name, or passing it as an ordinary value does not force it.
An implementation initializer is elaborated in the old $\Delta$: bindings are nonrecursive.

\newpage
\section{Typing and elaboration: records, primitives, and method syntax}
\begin{mathpar}

\inferrule*[right=E-Int]{n\in\mathbb Z}{\E n\Int n}

\inferrule*[right=E-String]{s\text{ a string}}{\E s\Str s}

\inferrule*[right=E-Bool]{b\in\{\mathsf{true},\mathsf{false}\}}{\E b\Bool b}


\inferrule*[right=E-Record]{\E{e_j}{A_j}{c_j}\quad(j\in J)\\\ell_j\text{ distinct}}
 {\E{\{\ell_j=e_j\}_{j\in J}}{\{\ell_j:A_j\}_{j\in J}}{\{\ell_j=c_j\}_{j\in J}}}


\inferrule*[right=E-Project]{\E e{\{\ell_j:A_j\}_{j\in J}}c\\ k\in J}
 {\E{e.\ell_k}{A_k}{c.\ell_k}}


\inferrule*[right=E-BinOp]{\E{e_1}{\argop\odot}{c_1}\\
 \E{e_2}{\argop\odot}{c_2}}
 {\E{e_1\mathbin{\odot}e_2}{\resop\odot}{c_1\mathbin{\odot}c_2}}
\end{mathpar}

\subsection*{Receiver-style method selection}
For a receiver type $A$ and label $\ell$, define the finite set of \emph{types}
\[
\mathcal M_{\Delta,A,\ell}
=\bigl\{(P,B)\ \bigm|\ (d:P)\in\Delta,\;
 P=\{\ell_j:A_j\}_{j\in J},\;
 \exists k\in J.\;\ell_k=\ell\land A_k=A\to B\bigr\}.
\]
Repeated occurrences of the same pair contribute one set member.
Let $\operatorname{labels}(A)$ be the field labels when $A$ is a record, and $\varnothing$ otherwise.
\begin{mathpar}

\inferrule*[right=E-Method]{\E e A c\\\ell\notin\operatorname{labels}(A)\\
 \mathcal M_{\Delta,A,\ell}=\{(P,B)\}}
 {\E{e.\ell}{B}{(\hole P).\ell\;c}}
\end{mathpar}
Zero candidates reject the selection; two distinct candidate types reject it as ambiguous.
This rule determines the request type, not its implementation. Actual record projection
has priority. Explicit qualification $\need P.\ell(e)$ avoids this receiver-syntax restriction.

\subsection*{The running qualified function}
With $\TS[A]=\{\kw{toString}:A\to\Str\}$, the source definition
\begin{lstlisting}
def display[A where ToString[A]](x: A): String = x.toString
\end{lstlisting}
has type and unresolved body
\[
\begin{aligned}
\kw{display}&:\forall A.\TS[A]\Rightarrow A\to\Str,\\
c_{\kw{display}}&=\Lambda A.\ilam{d}{\TS[A]}{
 \lam{x}{A}{(\hole{\TS[A]}).\kw{toString}\;x}}.
\end{aligned}
\]
Under $(d:\TS[A])::\Delta$, evidence resolution may select $d$.
After successful resolution:
\[
 t_{\kw{display}}
 =\Lambda A.\lam{d}{\TS[A]}{\lam{x}{A}{d.\kw{toString}\;x}},
 \qquad
 t_{\kw{display}}:\forall A.\TS[A]\to A\to\Str.
\]
A concrete call first becomes
$c_{\kw{display}}[\Str]\{\hole{\TS[\Str]}\}\;s$.
The request is solved in the \emph{caller's} lexical environment, while requests in
an ordinary function body are solved in that body's definition environment.

\newpage
\section{Admissible evidence and the synthesizer}
Evidence constructions are a restricted subset of the unresolved core:
\[
\boxed{p::=d\mid p[A]\mid p\,p\mid p\{p\}.}
\]
The only leaves are references into $\Delta$. New literals, lambdas, projections,
ordinary references from $\Gamma$, and unresolved requests are not synthesis results.
They may occur inside a previously checked first-class value installed in $\Delta$.

\begin{mathpar}

\inferrule*[right=P-Instance]{d:A\in\Delta}{\Pj d A}


\inferrule*[right=P-Type-App]{\Pj p{\forall\alpha.A}\\\wf B}
 {\Pj{p[B]}{A[B/\alpha]}}


\inferrule*[right=P-App]{\Pj{p_1}{A\to B}\\\Pj{p_2}A}
 {\Pj{p_1\,p_2}B}


\inferrule*[right=P-Implicit-App]{\Pj{p_1}{A\Rightarrow B}\\\Pj{p_2}A}
 {\Pj{p_1\{p_2\}}B}
\end{mathpar}
These rules \emph{check} a proposed construction. They do not prescribe a search strategy.
A pluggable algorithm supplies a partial judgment
\[
\Delta\vdash_{\mathcal S} A\rightsquigarrow p.
\]
The resolver rechecks $\Pj p A$ before lowering. The source arrow kind is checked before
erasure, even if two candidate types would erase to the same ordinary function type.

\subsection*{Exact-lookup placeholder}
Let $\Delta=[d_0:A_0,\ldots,d_{m-1}:A_{m-1}]$, nearest binding first.
\begin{mathpar}

\inferrule*[right=S-Exact]{0\le j<m\\A_j=A\\\forall k<j.\ A_k\ne A}
 {\Delta\vdash_{\mathrm{naive}}A\rightsquigarrow d_j}
\end{mathpar}
If no $j$ satisfies the premises, the placeholder fails. It does not instantiate a
universal type, apply a producer, recurse, or search for a hidden type argument.
Failure is not an uninhabitability result.

\subsection*{First-class polymorphic implementations}
A generic implementation may have the ordinary core type
\[
 i:\forall A,B.\kw{Repr}[A,B]\Rightarrow\TS[B]\Rightarrow\TS[A],
 \qquad\kw{Repr}[A,B]=\{\kw{repr}:A\to B\}.
\]
Given $r:\kw{Repr}[\kw{Person},\Str]$ and $s:\TS[\Str]$ in $\Delta$, the construction
\[
 p=i[\kw{Person}][\Str]\{r\}\{s\}
\]
satisfies $\Pj p{\TS[\kw{Person}]}$. It is permitted as the result of a stronger synthesizer,
but the exact-lookup placeholder does not discover it. The source program can construct
and install it explicitly:
\begin{lstlisting}
impl person: ToString[Person] = via[Person, String]{personName}{string}
ToString[Person].toString({ name = "Ada" })
\end{lstlisting}
There is no special type of implementation values: $i,r,s$ are ordinary values,
and $\mathsf{impl}$ only changes which context binds them.

\newpage
\section{Resolution and lowering}
\textbf{Type erasure} (identity on variables and base types):
\[
\begin{aligned}
\erase{A\to B}&=\erase A\to\erase B,&
\erase{A\Rightarrow B}&=\erase A\to\erase B,\\
\erase{\forall\alpha.A}&=\forall\alpha.\erase A,&
\erase{\{\ell_j:A_j\}_{j\in J}}&=\{\ell_j:\erase{A_j}\}_{j\in J}.
\end{aligned}
\]
For evidence, $\erase d=d$, $\erase{p[A]}=\erase p[\erase A]$, and both evidence
application forms erase to ordinary application. Ordinary and instance names are
embedded in one target namespace; freshness prevents capture.

\begin{mathpar}

\inferrule*[right=R-Var]{x:A\in\Gamma}{\R x x}

\inferrule*[right=R-Instance]{d:A\in\Delta}{\R d d}

\inferrule*[right=R-Need]{\Delta\vdash_{\mathcal S}A\rightsquigarrow p\\\Pj p A}
 {\R{\hole A}{\erase p}}


\inferrule*[right=R-Lambda]{\Gamma,x:A;\Delta\vdash c\Longrightarrow t}
 {\R{\lam{x}{A}{c}}{\lam{x}{\erase A}{t}}}

\inferrule*[right=R-Implicit-Lambda]{\Gamma;(d:A)::\Delta\vdash c\Longrightarrow t}
 {\R{\ilam{d}{A}{c}}{\lam{d}{\erase A}{t}}}


\inferrule*[right=R-App]{\R{c_1}{t_1}\\\R{c_2}{t_2}}
 {\R{c_1\,c_2}{t_1\,t_2}}

\inferrule*[right=R-Implicit-App]{\R{c_1}{t_1}\\\R{c_2}{t_2}}
 {\R{c_1\{c_2\}}{t_1\,t_2}}


\inferrule*[right=R-Type-Lambda]{\Gamma,\alpha;\Delta\vdash c\Longrightarrow t}
 {\R{\Lambda\alpha.c}{\Lambda\alpha.t}}

\inferrule*[right=R-Type-App]{\R c t}{\R{c[A]}{t[\erase A]}}


\inferrule*[right=R-Let]{\R{c_1}{t_1}\\\Gamma,x:A;\Delta\vdash c_2\Longrightarrow t_2}
 {\R{\letin{x:A}{c_1}{c_2}}{\letin{x}{t_1}{t_2}}}


\inferrule*[right=R-Impl]{\R{c_1}{t_1}\\\Gamma;(d:A)::\Delta\vdash c_2\Longrightarrow t_2}
 {\R{\impl{d}{A}{c_1}{c_2}}{\letin{d}{t_1}{t_2}}}


\inferrule*[right=R-Record]{\R{c_j}{t_j}\quad(j\in J)}
 {\R{\{\ell_j=c_j\}_{j\in J}}{\{\ell_j=t_j\}_{j\in J}}}

\inferrule*[right=R-Project]{\R c t}{\R{c.\ell}{t.\ell}}


\inferrule*[right=R-Int]{ }{\R n n}

\inferrule*[right=R-String]{ }{\R s s}

\inferrule*[right=R-Bool]{ }{\R b b}

\inferrule*[right=R-BinOp]{\R{c_1}{t_1}\\\R{c_2}{t_2}}
 {\R{c_1\mathbin{\odot}c_2}{t_1\mathbin{\odot}t_2}}
\end{mathpar}

\newpage
\section{Explicit target and independent typing}
\[
\begin{array}{rcl}
T,U&::=&\alpha\mid\Int\mid\Bool\mid\Str\mid T\to U\mid\forall\alpha.T\mid\{\ell_j:T_j\}_{j\in J}\\
t&::=&z\mid n\mid b\mid s\mid\lam z T t\mid t\,t\mid\Lambda\alpha.t\mid t[T]\\
&&\mid\letin z t t\mid\{\ell_j=t_j\}_{j\in J}\mid t.\ell
 \mid t_1\mathbin{\odot}t_2\\
H&::=&\cdot\mid H,\alpha\mid H,z:T
\end{array}
\]
Target type formation is Section~1 with only the ordinary arrow. All annotations and
type arguments must be well formed in $H$. All record labels are distinct.
\begin{mathpar}

\inferrule*[right=T-Var]{z:T\in H}{H\vdash z:T}

\inferrule*[right=T-Lambda]{H,z:T\vdash t:U}{H\vdash\lam z T t:T\to U}

\inferrule*[right=T-App]{H\vdash t_1:T\to U\\H\vdash t_2:T}{H\vdash t_1\,t_2:U}


\inferrule*[right=T-Type-Lambda]{H,\alpha\vdash t:T\\\alpha\notin\mathrm{FV}(H)}
 {H\vdash\Lambda\alpha.t:\forall\alpha.T}

\inferrule*[right=T-Type-App]{H\vdash t:\forall\alpha.T\\H\vdash U\;\kw{type}}
 {H\vdash t[U]:T[U/\alpha]}


\inferrule*[right=T-Let]{H\vdash t_1:T\\H,z:T\vdash t_2:U}
 {H\vdash\letin z{t_1}{t_2}:U}


\inferrule*[right=T-Record]{H\vdash t_j:T_j\quad(j\in J)}
 {H\vdash\{\ell_j=t_j\}_{j\in J}:\{\ell_j:T_j\}_{j\in J}}

\inferrule*[right=T-Project]{H\vdash t:\{\ell_j:T_j\}_{j\in J}\\k\in J}
 {H\vdash t.\ell_k:T_k}


\inferrule*[right=T-Int]{n\in\mathbb Z}{H\vdash n:\Int}

\inferrule*[right=T-String]{s\text{ a string}}{H\vdash s:\Str}

\inferrule*[right=T-Bool]{b\in\{\mathsf{true},\mathsf{false}\}}{H\vdash b:\Bool}

\inferrule*[right=T-BinOp]{H\vdash t_1:\argop\odot\\H\vdash t_2:\argop\odot}
 {H\vdash t_1\mathbin{\odot}t_2:\resop\odot}
\end{mathpar}

\subsection*{Phase invariant and index representation}
Let $H=\erase\Gamma\cup\erase\Delta$ merge the two term environments using the
same globally fresh names. The phase invariant is
\[
\boxed{\E e A c\quad\land\quad\R c t\quad\Longrightarrow\quad H\vdash t:\erase A.}
\]
The reference implementation independently checks the final judgment; no mechanized
metatheory is asserted here. The synthesizer is not a trusted source of typing derivations.

In the implementation, ordinary and instance references have independent de Bruijn indices.
The resolver maintains one embedding for each into the current target environment.
Every new term binder increments all previously embedded target indices. An ordinary
binder then extends only the $\Gamma$ embedding; an implementation binder extends only
the $\Delta$ embedding. A type binder weakens type indices in both environments and
in embedded evidence. Type substitution is capture avoiding.

\newpage
\section{Call-by-value evaluation}
\[
\begin{array}{rcl}
v&::=& n\mid b\mid s\mid\{\ell_j=v_j\}_{j\in J}\mid
 \langle\lambda z.t,\rho\rangle\mid\langle\Lambda\alpha.t,\rho\rangle,\\
\rho&::=&\cdot\mid\rho[z\mapsto v].
\end{array}
\]
Evaluation ignores type annotations and type arguments. Type abstraction creates a delayed
closure; type application forces it. Records are evaluated in canonical label order.

\begin{mathpar}

\inferrule*[right=V-Var]{\rho(z)=v}{\eval\rho z v}

\inferrule*[right=V-Int]{ }{\eval\rho n n}

\inferrule*[right=V-String]{ }{\eval\rho s s}

\inferrule*[right=V-Bool]{ }{\eval\rho b b}


\inferrule*[right=V-Lambda]{ }
 {\eval\rho{\lam z T t}{\langle\lambda z.t,\rho\rangle}}

\inferrule*[right=V-Type-Lambda]{ }
 {\eval\rho{\Lambda\alpha.t}{\langle\Lambda\alpha.t,\rho\rangle}}


\inferrule*[right=V-App]{
 \eval\rho{t_1}{\langle\lambda z.t,\rho_0\rangle}\\
 \eval\rho{t_2}{v_2}\\\eval{\rho_0[z\mapsto v_2]}{t}{v}}
 {\eval\rho{t_1\,t_2}{v}}


\inferrule*[right=V-Type-App]{\eval\rho{t_0}{\langle\Lambda\alpha.t,\rho_0\rangle}\\
 \eval{\rho_0}{t}{v}}
 {\eval\rho{t_0[T]}{v}}


\inferrule*[right=V-Let]{\eval\rho{t_1}{v_1}\\\eval{\rho[z\mapsto v_1]}{t_2}{v}}
 {\eval\rho{\letin z{t_1}{t_2}}{v}}


\inferrule*[right=V-Record]{\eval\rho{t_j}{v_j}\quad(j\in J)}
 {\eval\rho{\{\ell_j=t_j\}_{j\in J}}{\{\ell_j=v_j\}_{j\in J}}}

\inferrule*[right=V-Project]{\eval\rho t{\{\ell_j=v_j\}_{j\in J}}\\k\in J}
 {\eval\rho{t.\ell_k}{v_k}}


\inferrule*[right=V-BinOp]{\eval\rho{t_1}{v_1}\\\eval\rho{t_2}{v_2}\\
 \delta_{\odot}(v_1,v_2)=v}
 {\eval\rho{t_1\mathbin{\odot}t_2}v}
\end{mathpar}

\subsection*{Closed-program execution}
\begin{mathpar}

\inferrule*[right=Execute]{
 \cdot;[\,]\vdash e:A\rightsquigarrow c\\
 \cdot;[\,]\vdash c\Longrightarrow t\\
 \cdot\vdash t:\erase A\\\eval\cdot t v}
 {e\Downarrow_{\mathcal S}v}
\end{mathpar}
There is no runtime evidence search. Resolving an implementation binding produces an
ordinary target \textsf{let}, so its value can close over ordinary runtime variables.
The compiled value remains first-class; registration affects only lexical evidence scope.

\end{document}
